\section{The world drift and an honest error bar}
\label{sec:cond-drift}

\paragraph{Why world-fixed.} With the thrusters dead the boat moves only with the
water and the wind, both fixed in the \emph{world} frame (up to a heading effect
treated in Section~\ref{sec:cond-observe}). Measuring the drift in the world
frame needs no heading at all: GPS velocity is already east and north. A heading
is needed when residual velocities are rotated into the body.

\subsection{The pooled estimate}

Each window $w$ gives a mean $\bm m_w$ whose two components have variance
$\Szero/T_w$ (Section~\ref{sec:cond-noise}). Longer windows are more precise, so
they should count more. The minimum-variance way to combine independent
estimates of one quantity is to weight each by its inverse variance. Minimizing
\[
  \chi^2(\bm d) = \sum_w \frac{T_w\,\|\bm m_w - \bm d\|^2}{\Szero}
\]
over $\bm d$ (set the gradient $-2\sum_w T_w(\bm m_w - \bm d)/\Szero$ to zero) gives
\[
  \dhat = \frac{\sum_w T_w\,\bm m_w}{\sum_w T_w},\qquad
  \Var(\hat d_j) = \frac{\Szero}{\sum_w T_w}\quad\text{per component } j.
\]
It is the duration-weighted mean, and its variance is that of one long window of
the total neutral time, which is what one expects if the windows are independent.

\paragraph{Overdispersion.} The model says the window means scatter about
$\dhat$ with variance $\Szero/T_w$. If they scatter more, the model is missing
something: a gust that lasted longer than a window, weathervaning, a slow swing
of the current. The check is the minimized $\chi^2$ divided by its degrees of
freedom,
\[
  D_{\mathrm{od}} = \frac{\chi^2(\dhat)}{\mathrm{df}},\qquad \mathrm{df} = 2W - p,
\]
two components per window, less the $p$ fitted parameters ($p = 2$ for a constant
drift, $4$ with the windage term of Section~\ref{sec:cond-observe}). Under the model
$\E[D_{\mathrm{od}}] = 1$. When $D_{\mathrm{od}} > 1$ the error bar is scaled up by $\sqrt{D_{\mathrm{od}}}$; it is never
scaled down, because a $D_{\mathrm{od}}$ below one from a handful of windows is luck, not
evidence of extra precision.

\paragraph{Correlation between windows.} $D_{\mathrm{od}}$ catches scatter, but not
\emph{memory}. Noise that wanders on minute scales makes neighbouring window means
alike: they scatter little about each other, so $D_{\mathrm{od}}$ stays near one, yet the pooled
mean is far less certain than $W$ independent windows would make it. The
factor for window means with lag-$\ell$ correlation $r^\ell$, summed as a geometric series,
\[
  \Var\Big(\frac1W\sum_w x_w\Big) = \frac{\sigma^2}{W}\Big(1 + 2\sum_{\ell=1}^{W-1}\big(1-\tfrac{\ell}{W}\big)r^\ell\Big)
  \;\approx\; \frac{\sigma^2}{W}\Big(1 + \frac{2r}{1-r}\Big) = \frac{\sigma^2}{W}\,\frac{1+r}{1-r}.
\]
The estimator measures $r$ as the lag-1 correlation of the window-mean residuals,
adds $1/W$ (the same demeaning bias as before, in its large-sample form), and clamps
it to $[0, 0.6]$ so that three windows cannot produce an absurd factor. Putting
the three pieces together,
\[
  \boxed{\ \sed^2 = \frac{\Szero}{\sum_w T_w}\cdot\max(1, D_{\mathrm{od}})\cdot\frac{1+r}{1-r}\ }
\]
per component. Figure~\ref{fig:cond-coverage} shows why the last factor is there:
without it, the 95\,\% interval covers the truth only 88\,\% of the time when the
noise has a 6\,s memory, which is exactly the regime where $\Szero$ itself reads
low because of the cap on $\tau_c$.

\begin{figure}[ht]
\centering
\begin{tikzpicture}
\begin{axis}[cond axis, width=0.6\linewidth, height=5cm, ybar, bar width=9pt,
  xlabel={noise correlation time $\tau_c$ (s)}, ylabel={coverage of 95\,\% interval},
  symbolic x coords={1,3,6}, xtick=data, ymin=0.8, ymax=1.0,
  legend columns=2, legend style={at={(0.5,1.03)}, anchor=south},
  grid=major, grid style={gray!20}]
  \addplot[fill=red!40, draw=red] table[x=tau, y=plain] {figures/conditions/coverage.dat};
  \addlegendentry{$\Szero$ and $D_{\mathrm{od}}$ only}
  \addplot[fill=blue!40, draw=sA] table[x=tau, y=inflated] {figures/conditions/coverage.dat};
  \addlegendentry{with $(1+r)/(1-r)$}
  \draw[black, dashed] (rel axis cs:0,0.75) -- (rel axis cs:1,0.75);
\end{axis}
\end{tikzpicture}
\caption{How often $|\hat d - d| \le 1.96\,\sed$ on 400 synthetic sessions of
12 windows (96\,s), AR(1) noise with the $\tau_c$ shown, $\tau_c$ capped at one
batch as in production. The dashed line is the nominal 95\,\%. A simplified
re-implementation for illustration; the production code's own coverage tests
(300 runs of 300\,s) give 0.97, 0.96 and 0.94 at 1, 3 and 6\,s. Data:
\code{figures/conditions/coverage.dat}.}
\label{fig:cond-coverage}
\end{figure}

\subsection{What the drift contains}

The pooled $\dhat$ is neutral velocity over ground: water current, wind
leeway and receiver bias are inseparable in this estimate. Characterize subtracts
its local pre-hold neutral baseline rather than this pooled drift.

\file{mhseals\_testing/analysis/conditions.py (\_pool, world\_drift)}
